\documentclass{article}
\usepackage{amsmath, amssymb, amsthm}
\usepackage{hyperref}
\usepackage{geometry}
\geometry{margin=1in}

\title{Erd\H{o}s Problem \#959}
\date{Last updated: 2025-08-31}
\author{Source: erdosproblems.com}

\begin{document}
\maketitle

\noindent\textbf{Status:} OPEN \quad \textbf{No prize}

\noindent\textbf{Tags:} geometry, distances

\medskip
\noindent\textbf{Problem Statement:}

Let $A\subset \mathbb{R}^2$ be a set of size $n$ and let $\{d_1,\ldots,d_k\}$ be the set of distinct distances determined by $A$. Let $f(d)$ be the number of times the distance $d$ is determined, and suppose the $d_i$ are ordered such that\[f(d_1)\geq f(d_2)\geq \cdots \geq f(d_k).\]Estimate\[\max (f(d_1)-f(d_2)),\]where the maximum is taken over all $A$ of size $n$.

\bigskip
\noindent\textbf{Remarks:}

More generally, one can ask about\[\max (f(d_r)-f(d_{r+1})).\]Clemen, Dumitrescu, and Liu \cite{CDL25}, have shown that\[\max (f(d_1)-f(d_2))\gg n\log n.\]More generally, for any $1\leq k\leq \log n$, there exists a set $A$ of $n$ points such that\[f(d_r)-f(d_{r+1})\gg \frac{n\log n}{r}.\]They conjecture that $n\log n$ can be improved to $n^{1+c/\log\log n}$ for some constant $c>0$.

\bigskip
\noindent\textbf{References:}

[CDL25] F. Clemen, A. Dumitrescu, and D. Liu, On multiplicities of interpoint distances. arXiv:2505.04283 (2025).

\bigskip
\noindent\small{Source: \url{https://www.erdosproblems.com/959}}
\end{document}
